% GENERATED FILE. Edit canonical lessons, then run npm run generate:books.
% canonical-source-hash: fnv1a64:193506f5f35db3cf
% canonical-lessons: TE-C188-prasadam, TE-C188-ariselu, TE-C188-toranam, TE-C188-bottu, TE-C188-gorintaku

\chapter{Prasadam, Rice-flour sweets, Festoon of mango leaves hung over a doorway, Bindi, Henna}
\label{ch:te-prasadam-rice-flour-sweets-festoon-of-mango-leaves-hung-over-a-doorway-bindi-henna}
\hlchaptermodality{188}

\begin{chapteropening}
\textbf{By the end of this chapter:} \emph{I can say prasadam, rice-flour sweets, a festoon of mango leaves hung over a doorway, a bindi and henna.}

Complete the last lesson of chapter 188: I can say prasadam, rice-flour sweets, a festoon of mango leaves hung over a doorway, a bindi and henna.
\end{chapteropening}

\section[\texorpdfstring{\te{ప్రసాదం} (prasādaṁ)}{ప్రసాదం (prasādaṁ)}]{\te{ప్రసాదం} (prasādaṁ) — prasadam, food offered to a deity and shared}

\label{lesson:TE-C188-prasadam}



\begin{quote}
Before the new one: say the Telugu for a sparrow, then the Telugu for a koel.
\end{quote}

\subsection*{You'll want to know: \te{ప్రసాదం}}
\textbf{\te{ప్రసాదం}} — \emph{prasādaṁ} — \textquotedblleft{}prasadam, food offered to a deity and shared\textquotedblright{}.

\begin{grammarlens}[title={What you've built}]
One more word for this chapter.
\end{grammarlens}

\subsection*{Your turn}
\begin{itemize}
\raggedright
  \item \emph{Say it:} \emph{prasādaṁ}
  \item \emph{Say it:} \emph{prasādaṁ}, once more
  \item \emph{Say it:} say \emph{prasādaṁ}
  \item \emph{Recall:} say the Telugu for a sparrow, then the Telugu for a koel, then say \emph{prasādaṁ} again
\end{itemize}

\begin{tcolorbox}[breakable,colback=teal!4,colframe=teal!35!black,title={Before you move on}]
What does \te{ప్రసాదం} mean? (\textquotedblleft{}Prasadam, food offered to a deity and shared\textquotedblright{}.) Say it once more.
\end{tcolorbox}
\section[\texorpdfstring{\te{అరిసెలు} (ariselu)}{అరిసెలు (ariselu)}]{\te{అరిసెలు} (ariselu) — rice-flour sweets; made with jaggery at Sankranti}

\label{lesson:TE-C188-ariselu}



\begin{quote}
Before the new one: say the Telugu for a koel, then the Telugu for prasadam.
\end{quote}

\subsection*{You'll want to know: \te{అరిసెలు}}
\textbf{\te{అరిసెలు}} — \emph{ariselu} — \textquotedblleft{}rice-flour sweets; made with jaggery at Sankranti\textquotedblright{}.

\begin{grammarlens}[title={What you've built}]
One more word for this chapter.
\end{grammarlens}

\subsection*{Your turn}
\begin{itemize}
\raggedright
  \item \emph{Say it:} \emph{ariselu}
  \item \emph{Say it:} \emph{ariselu}, once more
  \item \emph{Say it:} say \emph{ariselu}
  \item \emph{Recall:} say the Telugu for a koel, then the Telugu for prasadam, then say \emph{ariselu} again
\end{itemize}

\begin{tcolorbox}[breakable,colback=teal!4,colframe=teal!35!black,title={Before you move on}]
What does \te{అరిసెలు} mean? (\textquotedblleft{}Rice-flour sweets; made with jaggery at Sankranti\textquotedblright{}.) Say it once more.
\end{tcolorbox}
\section[\texorpdfstring{\te{తోరణం} (tōraṇaṁ)}{తోరణం (tōraṇaṁ)}]{\te{తోరణం} (tōraṇaṁ) — a festoon of mango leaves hung over a doorway}

\label{lesson:TE-C188-toranam}



\begin{quote}
Before the new one: say the Telugu for prasadam, then the Telugu for rice-flour sweets.
\end{quote}

\subsection*{You'll want to know: \te{తోరణం}}
\textbf{\te{తోరణం}} — \emph{tōraṇaṁ} — \textquotedblleft{}a festoon of mango leaves hung over a doorway\textquotedblright{}.

\begin{grammarlens}[title={What you've built}]
One more word for this chapter.
\end{grammarlens}

\subsection*{Your turn}
\begin{itemize}
\raggedright
  \item \emph{Say it:} \emph{tōraṇaṁ}
  \item \emph{Say it:} \emph{tōraṇaṁ}, once more
  \item \emph{Say it:} say \emph{tōraṇaṁ}
  \item \emph{Recall:} say the Telugu for prasadam, then the Telugu for rice-flour sweets, then say \emph{tōraṇaṁ} again
\end{itemize}

\begin{tcolorbox}[breakable,colback=teal!4,colframe=teal!35!black,title={Before you move on}]
What does \te{తోరణం} mean? (\textquotedblleft{}A festoon of mango leaves hung over a doorway\textquotedblright{}.) Say it once more.
\end{tcolorbox}
\section[\texorpdfstring{\te{బొట్టు} (boṭṭu)}{బొట్టు (boṭṭu)}]{\te{బొట్టు} (boṭṭu) — a bindi, a dot worn on the forehead}

\label{lesson:TE-C188-bottu}



\begin{quote}
Before the new one: say the Telugu for rice-flour sweets, then the Telugu for a festoon of mango leaves hung over a doorway.
\end{quote}

\subsection*{You'll want to know: \te{బొట్టు}}
\textbf{\te{బొట్టు}} — \emph{boṭṭu} — \textquotedblleft{}a bindi, a dot worn on the forehead\textquotedblright{}.

\begin{grammarlens}[title={What you've built}]
One more word for this chapter.
\end{grammarlens}

\subsection*{Your turn}
\begin{itemize}
\raggedright
  \item \emph{Say it:} \emph{boṭṭu}
  \item \emph{Say it:} \emph{boṭṭu}, once more
  \item \emph{Say it:} say \emph{boṭṭu}
  \item \emph{Recall:} say the Telugu for rice-flour sweets, then the Telugu for a festoon of mango leaves hung over a doorway, then say \emph{boṭṭu} again
\end{itemize}

\begin{tcolorbox}[breakable,colback=teal!4,colframe=teal!35!black,title={Before you move on}]
What does \te{బొట్టు} mean? (\textquotedblleft{}A bindi, a dot worn on the forehead\textquotedblright{}.) Say it once more.
\end{tcolorbox}
\section[\texorpdfstring{\te{గోరింటాకు} (gōriṇṭāku)}{గోరింటాకు (gōriṇṭāku)}]{\te{గోరింటాకు} (gōriṇṭāku) — henna}

\label{lesson:TE-C188-gorintaku}



\begin{quote}
Before the new one: say the Telugu for a festoon of mango leaves hung over a doorway, then the Telugu for a bindi.
\end{quote}

\subsection*{You'll want to know: \te{గోరింటాకు}}
\textbf{\te{గోరింటాకు}} — \emph{gōriṇṭāku} — \textquotedblleft{}henna\textquotedblright{}.

\begin{grammarlens}[title={What you've built}]
One more word for this chapter.
\end{grammarlens}

\subsection*{Your turn}
\begin{itemize}
\raggedright
  \item \emph{Say it:} \emph{gōriṇṭāku}
  \item \emph{Say it:} \emph{gōriṇṭāku}, once more
  \item \emph{Say it:} say \emph{gōriṇṭāku}
  \item \emph{Recall:} say the Telugu for a festoon of mango leaves hung over a doorway, then the Telugu for a bindi, then say \emph{gōriṇṭāku} again
\end{itemize}

\begin{tcolorbox}[breakable,colback=teal!4,colframe=teal!35!black,title={Before you move on}]
What does \te{గోరింటాకు} mean? (\textquotedblleft{}Henna\textquotedblright{}.) Say it once more.
\end{tcolorbox}
